\documentclass[10pt,a4paper]{amsart}
\usepackage[margin=19mm]{geometry}
\usepackage{amsmath,amssymb,mathtools}
\usepackage[scr=rsfs,cal=euler]{mathalfa}
\usepackage{xfrac}
\usepackage[unicode,hidelinks,bookmarks=false]{hyperref}
\newcommand{\ie}{i.e.\ }
\newcommand{\cf}{cf.\ }
\newcommand{\resp}{respectively\ }
\newcommand{\Subsection}{\subsection}
\providecommand{\nobreakdash}{\nobreak\hbox{-}}
\setlength{\emergencystretch}{2em}
\multlinegap0pt
\usepackage{xcolor}
\usepackage{braket}
\usepackage{booktabs}
\usepackage{multirow}
\usepackage{amssymb}
\usepackage{nicefrac}
\usepackage{mathtools}
\usepackage{enumerate}
\usepackage{bm}
\usepackage{float}
\usepackage{dsfont}
\usepackage{algorithm}\usepackage{algpseudocode}
\newtheorem{theorem}{Theorem}
\renewcommand{\H}{{\scriptscriptstyle\mathsf{H}}}
\def\bD{{\mathbf{D}}}
\def\bP{{\mathbf{P}}}
\def\bR{{\mathbf{R}}}
\def\nn{\nonumber}
\title{On the Ergodic Capacity for SIM-Aided Holographic MIMO Communications}
\author{Anastasios Papazafeiropoulos, Ioannis Bartsiokas, Dimitra I. Kaklamani, Iakovos S. Venieris}
\date{Source: arXiv:2512.22068v1 (CC BY 4.0)}
\begin{document}
\maketitle

\noindent\emph{Source and licence.} On the Ergodic Capacity for SIM-Aided Holographic MIMO Communications. Version arXiv:2512.22068v1 (CC BY 4.0), distributed by arXiv under the Creative Commons Attribution 4.0 International licence, http://creativecommons.org/licenses/by/4.0/, as recorded in the arXiv licence metadata for this exact version and shown on the abstract page https://arxiv.org/abs/2512.22068v1. Full licence notice: \texttt{licenses/arxiv-2512.22068-CC-BY-4.0.txt}.\par\smallskip

\noindent\textit{Editorial scope.} Counted unit: the article's theorem th1 (source0.tex lines 575--582). The article's complete original proof of it is delivered with it (source0.tex lines 583--585, one \texttt{proof} environment of 3 lines). No mathematics is written, repaired or re-derived: the statement and the proof are the article's own lines, in the article's own notation, and no retained line is substituted or re-flowed.  Retained uncounted as the source-local context the proof needs: lines 820--820.  Nothing was omitted from any retained span, so the package carries no omission record.  No retained line contains a citation, so the wrapper carries no bibliography and no bibliography entry.  No retained line contains a figure environment, an image-inclusion command or drawing code, so no image is delivered and the package declares no asset dependency.  Editorial formatting only, in addition: an AMS \texttt{amsart} preamble that restates the article's own declarations and macros used by the retained lines, namely the environments \texttt{theorem} on separate counters, together with the standard packages those lines need.  The retained environments print in this document as Theorem 1 in order of appearance, whereas the article prints the same items with the numbering of its own full document.  The complete native source of the article as distributed in the arXiv e-print for this exact version is bundled unchanged as \texttt{sources/191713/source0.tex}.
		\section{Proof of Theorem~\ref{th1}}\label{th1proof}	

	\begin{theorem}\label{th1}
	The ergodic capacity of a  SIM-aided HMIMO system admits a lower bound given by
		\begin{align}
			C_{\mathrm{erg}}&\ge 	C_{\mathrm{LB}}=N_{t} \log_{2}\bigg(\!\!1+\frac{\rho}{N_{t}} \exp \Big(\frac{1}{N_{t}}\big(\textcolor{black}{\sum_{i=0}^{s-1}\psi(t-i)}\nn\\
			&+\ln |\bP\bP^{\H}|+\ln |\bD^{\H}\bD|+\ln |\bR_{\mathrm{T}}|+|\bR_{\mathrm{R}}|\big)\!\Big)\!\bigg)\label{capacity0}
		\end{align}
		with $\psi(x)$ being the digamma function [13, eq. (8.360.1)].
	\end{theorem}

	\begin{proof}
		Please see Appendix~\ref{th1proof}.	
	\end{proof}

\end{document}
